\documentclass[12pt]{article}
\usepackage{amsmath,amssymb}

\begin{document}
\section*{{2019 AMC 12B Problems/Problem 12}}
Source: AoPS Wiki

\subsection*{Solution}
First, note by the Pythagorean Theorem in $\triangle ABC$ that $AC = \sqrt{2}$ . Now, the equal perimeter condition means that $BC + BA = 2 = CD + DA$ , since side $AC$ is common to both triangles and thus can be discounted. This relationship, in combination with the Pythagorean Theorem in $\triangle ACD$ , gives $AC^2+CD^2=\left(\sqrt{2}\right)^2+\left(2-DA\right)^2=DA^2$ . Hence $2 + 4 - 4DA + DA^2 = DA^2$ , so $DA = \frac{3}{2}$ , and thus $CD = \frac{1}{2}$ . Next, since $\angle BAC = 45^{\circ}$ , $\sin{\left(\angle BAC\right)} = \cos{\left(\angle BAC\right)} = \frac{1}{\sqrt{2}}$ . Using the lengths found above, $\sin{\left(\angle CAD\right)} = \frac{\left(\frac{1}{2}\right)}{\left(\frac{3}{2}\right)} = \frac{1}{3}$ , and $\cos{\left(\angle CAD\right)} = \frac{\sqrt{2}}{\left(\frac{3}{2}\right)} = \frac{2 \sqrt{2}}{3}$ . Thus, by the addition formulae for $\sin$ and $\cos$ , we have \[\begin{split}\sin{\left(\angle BAD\right)}&=\sin{\left(\angle BAC + \angle CAD\right)}\\&=\sin{\left(\angle BAC\right)}\cos{\left(\angle CAD\right)}+\cos{\left(\angle BAC\right)}\sin{\left(\angle CAD\right)}\\&=\frac{1}{\sqrt{2}}\cdot\frac{2 \sqrt{2}}{3} + \frac{1}{\sqrt{2}}\cdot\frac{1}{3} \\ &= \frac{2 \sqrt{2} + 1}{3 \sqrt{2}}\end{split}\] and \[\begin{split}\cos{\left(\angle BAD\right)}&=\cos{\left(\angle BAC + \angle CAD\right)}\\&=\cos{\left(\angle BAC\right)}\cos{\left(\angle CAD\right)}-\sin{\left(\angle BAC\right)}\sin{\left(\angle CAD\right)}\\&=\frac{1}{\sqrt{2}}\cdot\frac{2 \sqrt{2}}{3} - \frac{1}{\sqrt{2}}\cdot\frac{1}{3} \\ &= \frac{2 \sqrt{2} - 1}{3 \sqrt{2}}\end{split}\] Hence, by the double angle formula for $\sin$ , $\sin{\left(2\angle BAD\right)} = 2\sin{\left(\angle BAD\right)}\cos{\left(\angle BAD\right)} = \frac{2(8-1)}{18} = \boxed{\textbf{(D) } \frac{7}{9}}$ .

\end{document}
